\section{Introduction}

Choosing a route and choosing a \emph{routing policy} are different decisions.
A navigation service may minimise distance, or estimated travel time, or may
spread demand across alternatives, or may avoid corridors whose travel time is
unpredictable. Each of these is an established objective with its own
literature. Production systems commit to one and apply it uniformly.

Whether that commitment costs anything depends on the traffic, not on the
planner. A distance-minimising policy is unimprovable when the shortest
corridor is uncongested and indefensible when it is saturated. A policy that
reacts to measured travel times is useful when congestion builds slowly and
counterproductive when a large guided cohort acts together on information that
is already stale~\cite{arnott1991information,acemoglu2018braess}. Whether these
regimes exist, where their boundaries lie, and whether they can be recognised
before a policy is activated, are empirical questions.

This paper answers them for a controlled microsimulation environment. The
question is:

\begin{quote}
\emph{Given a traffic context observable before a routing policy is activated,
which established policy should be activated, why is it the right one in that
regime, and when should the selector decline to adapt at all?}
\end{quote}

\paragraph{Research questions.}
\textbf{RQ1} Under which traffic, signal, information, penetration, capacity and
disruption conditions do established routing policies exhibit complementary
strengths?
\textbf{RQ2} Can the preference boundaries between them be explained by
mechanistic traffic variables rather than only by a statistical learner?
\textbf{RQ3} Can a machine-learning selector identify useful selection structure
beyond a strong mechanistic baseline?
\textbf{RQ4} When the operating context lies outside the training support, or
uncertainty is too high, can the system abstain safely and fall back to a fixed
policy?

\paragraph{Contributions.}
\textbf{(1)} A systematic characterisation of the operating regions in which
four established routing policies exhibit complementary performance, under
congestion, signal timing, information lag, guidance penetration, disruption
and alternative capacity.
\textbf{(2)} A mechanism-informed explanation of policy switching using
interpretable dimensionless traffic descriptors, so that the boundary is a
statement about traffic and not about a model.
\textbf{(3)} A decision-oriented context-aware selector, evaluated by decision
regret against both a strong fixed baseline and the mechanistic rule --- the
latter being the comparison that matters.
\textbf{(4)} A selective mechanism that recognises support and uncertainty
limits and abstains to a fixed policy, evaluated separately under five kinds of
distribution shift.

We introduce no routing algorithm, no learning algorithm and no multi-criteria
decision method. Every policy and every model class used here is established
and cited. The contribution is the characterisation, the mechanism, the
decision evaluation and the safety envelope.

\section{Related work}

\paragraph{Routing objectives.} Distance- and travel-time-minimising route
choice sits within the equilibrium tradition founded by
Wardrop~\cite{wardrop1952}. Load-balancing route guidance, in which vehicles
are distributed over a set of travel-time-acceptable alternatives according to
predicted link load, is established; we use the flow-balanced $k$-shortest-path
construction of Pan et al.~\cite{pan2013proactive}. Reliability-aware guidance,
in which a dispersion term is added to expected travel time, follows
Sen et al.~\cite{sen2001meanvariance}. Eco-routing, which minimises an
emissions proxy, is established~\cite{boriboonsomsin2012eco}; whether it is
\emph{distinct} from time-based routing in a given network is an empirical
question that we test rather than assume.

\paragraph{Information, penetration and instability.} That providing
information to drivers need not reduce congestion has been understood since
Arnott et al.~\cite{arnott1991information}, and has been given a sharp
game-theoretic form by Acemoglu et al.~\cite{acemoglu2018braess}. The benefit
of guidance depends on how much of the demand is guided~\cite{yang1998penetration}.
These results motivate treating penetration and information lag as
\emph{factors}, not as fixed implementation details.

\paragraph{Per-instance algorithm selection.} Choosing among algorithms by
instance features is Rice's problem~\cite{rice1976}, with a mature
methodology~\cite{xu2008satzilla,smithmiles2009,kotthoff2014survey} and
standard reference points --- the virtual best and single best solver --- from
the benchmarking literature~\cite{bischl2016aslib}. Instance-space analysis
asks where in a feature space each algorithm is
strong~\cite{smithmiles2014instance}, which is precisely the question RQ1 asks
of routing policies. Evaluating a predictor by the cost of the decision it
induces rather than by its prediction error is the predict-then-optimise
distinction of Elmachtoub and Grigas~\cite{elmachtoub2022spo}.

\paragraph{Abstention and distribution shift.} Deciding not to decide is
classical~\cite{chow1970reject}. Split conformal prediction supplies
distribution-free intervals~\cite{vovk2005alrw,lei2018conformal}, but its
coverage guarantee rests on exchangeability, which covariate shift
breaks~\cite{shimodaira2000,tibshirani2019covshift}. We therefore use a
conformal gate as a calibrated in-distribution instrument and a separate
support gate for extrapolation, and say so explicitly rather than claiming
that conformal prediction solves the out-of-distribution problem.

\paragraph{Positioning.} We do not claim that regime-dependent routing has
never been studied. Each component above has precedent. What is offered is the
combination: an explicitly established policy portfolio selected per context in
a physical microsimulation, evaluated by decision regret against both a best
fixed policy \emph{and} a mechanistic rule, with mechanism-informed features,
distribution shift reported separately by shift type, and a support-and-
confidence gate that abstains.

\section{Problem formulation}

A \emph{context} $s$ is a set of operating conditions known before a policy is
activated. A \emph{policy} $a \in \mathcal{P}$ maps individual vehicles to
routes by its own objective. The decision is the map $\pi: s \mapsto a$, taken
once per context.

\paragraph{Two levels.} At level~1 a context selects a policy; at level~2 the
selected policy routes vehicles. This is not route prediction, not
vehicle-level action learning, not signal control and not joint
routing--signal optimisation. The policies are fixed, published objectives;
only the choice among them is learned.

\paragraph{Decision cost and regret.} Writing $C(s,a)$ for the cost of running
policy $a$ in context $s$,
\begin{equation}
R(s,a) \;=\; C(s,a) \;-\; \min_{b \in \mathcal{P}} C(s,b).
\end{equation}
Two reference points from the algorithm-selection literature~\cite{bischl2016aslib}
bound what is achievable. The \textbf{single best policy} (SBS) is the one
fixed policy with the lowest mean cost, chosen on training contexts only; it is
what a deployed fixed-policy system achieves, and it is the baseline any
selector must beat to be worth deploying. The \textbf{virtual best policy}
(VBS) chooses the best policy for each context using that context's realised
outcomes. \emph{VBS is retrospective. It is not an algorithm, it cannot be
deployed, and it is used here only as an upper reference.} The distance between
them is the headroom available to any selector.

\paragraph{Criteria.} Three criteria are measured (Table~\ref{tab:criteria}):
mean journey time C1, \CO{} emissions C2, and total stopped delay C3 --- the
time vehicles spend stationary in queues. Policies are first compared by Pareto
dominance over all three. Scalarisation is used only for sensitivity, through
three declared normalised weight profiles. No monetary value of time, value of
reliability or carbon price is used: no verified source for them was available
in this study, and inventing one would make the conclusion an artefact of the
invented number.

\section{Experimental methodology}

\paragraph{Network.} Figure~\ref{fig:network} shows the environment. One
origin--destination pair is served by three corridors between a common diverge
and a common merge: a short, signalised, two-lane central arterial; a longer
and slower signalised street with one or two lanes; and a longest,
uninterrupted single-lane bypass. Four minor streets carry independent demand
through the same signals, so that a routing policy that loads a corridor also
delays traffic that is not being routed. Because every junction permits
straight-through movements only, the set of loopless paths from origin to
destination is exactly these three: the candidate path catalogue is complete,
not a $k$-shortest approximation.

Two structural choices are deliberate. The destination portal gives each
approach a dedicated exit lane, so that an alternative corridor is not
penalised by junction priority rather than by its own capacity; and the
diverge is not subject to a turn-radius speed ceiling, which in schematic node
geometry would cap the bypass at a value belonging to the drawing rather than
to the road. Both were identified by capacity validation before the experiment
was frozen.

\paragraph{Measured capacity.} Each corridor was driven to saturation alone and
its discharge measured. Queue-discharge saturation flow is
\NumSatArterial~veh/h/lane on the arterial and \NumSatNorth~veh/h/lane on the
slower street, consistent to within \NumSatSpread{} across green ratios; the
uninterrupted bypass carries \NumCapBypass~veh/h, below queue-discharge
saturation flow as free-flow headways require. Capacity is therefore a measured
property of this network, not a textbook constant asserted about it, and the
demand levels of Table~\ref{tab:factors} follow from it.

\paragraph{Policies.} The four policies are given in
Table~\ref{tab:portfolio} and Figure~\ref{fig:portfolio}. All read traffic
state from one shared observation buffer: a policy acting at time $t$ may read
link travel times sampled no later than $t-\Delta$, over a window of 120\,s,
and falls back to free-flow values before any admissible sample exists. No
policy sees the future and none sees any realised outcome of the run it is part
of. Policy parameters were frozen before the experiment: admissibility
tolerance $\varepsilon=0.20$ and a 120\,s assignment window for P3,
$\lambda=1.0$ for P4.

Each policy was validated before use. Twelve known-answer tests check that the
observation lag hides newer samples, that the free-flow fallback engages, that
P1 ignores congestion, that P2 follows the observed mean, that P4 prefers the
steadier of two paths with equal means, and that P3 respects its admissibility
tolerance. A two-path integration test inside the simulator, with equal path
lengths and a 2:1 capacity ratio, confirms that P3 splits demand
\NumTwoPathRatio{} against the capacity ratio of 2.00 while P1 commits the
entire cohort to one path.

\paragraph{Factors and common random numbers.} Six factors
(Table~\ref{tab:factors}, Figure~\ref{fig:design}) are crossed fully, giving
\NumContexts{} contexts, each run with \NumSeeds{} seeds and all four policies.
The vehicle set, departure times, guided/unguided labelling, habitual corridor
choice, background traffic and the incident realisation are functions of the
seed and the context alone, never of the policy. Two runs of a context under
different policies therefore face identical vehicles, and every comparison in
this paper is paired.

Unguided drivers are not re-routed by anything: they follow a habitual corridor
drawn once from a frozen logit over free-flow path times. This is what makes
penetration meaningful --- at penetration $p$, a fraction $1-p$ of corridor
demand is inert with respect to the policy under test.

\paragraph{Disruption as risk, not as a known event.} When the incident regime
is active, one lane closure is drawn --- location, onset and duration --- from a
declared distribution seeded by the run seed. Within a context and seed the
realisation is identical across all four policies, preserving the pairing
exactly; across the seeds of a context it differs, so the context carries
genuine travel-time risk that no policy can know in advance. An incident that
occurred at an identical time and place in every replication would be a known
event rather than a risk, and could not test a reliability-aware policy at all.

\paragraph{Statistical treatment.} The experimental unit is the context, not
the vehicle. An advantage of $a$ over $b$ is \emph{resolved} only when, across
the paired seed differences, $|\overline{d}| > 2\,\mathrm{SE}(\overline{d})$.
Unresolved differences are treated as ties throughout, so that no claim in this
paper rests on a difference the seeds cannot separate.

\section{Mechanistic regime model}

A selector that fires on raw demand explains nothing. We therefore express each
context through dimensionless descriptors with physical meaning, all computable
before policy activation:
network saturation $x_1 = D/\sum_c \mathrm{cap}_c$;
saturation of the shortest corridor $x_2 = D/\mathrm{cap}_C$;
alternative capacity share $x_3$;
effective green ratio $x_4$;
penetration $x_5$;
information staleness $x_6 = \Delta / T^{\mathrm{ff}}_C$, the lag as a fraction
of the free-flow trip time it is meant to describe;
incident regime $x_7$; and
controlled-flow saturation $x_8 = p\,D/\mathrm{cap}_C$, the share of the
shortest corridor's capacity that the policy actually commands.

Corridor capacity uses the measured saturation flows, effective green and lane
count, so $x_1$, $x_2$ and $x_8$ are genuine degrees of saturation rather than
scaled demand. A route-overlap index is not used: the three corridors share only
their access and egress edges, so overlap is structurally constant here and
carries no information.

The mechanistic rule B1 is an axis-aligned region rule on $x_1$--$x_8$, of depth
at most two, fitted on training contexts only. Its thresholds are reported as
\emph{brackets between adjacent sampled levels}, never as exact continuous
values: the experiment samples a grid, and a grid cannot locate a boundary more
finely than its own spacing.

\section{Context-aware selector}

\paragraph{Target.} The selector predicts \emph{advantage}, not cost. For each
policy $a$,
\begin{equation}
\Delta_a(x) \;=\; C_{\mathrm{SBS}}(x) \;-\; C_a(x),
\end{equation}
where the single best policy is determined on \emph{training contexts only}.
A positive $\Delta_a$ means $a$ is better in that context than the fixed policy
a deployed system would otherwise run, so the quantity being learned is exactly
the quantity a decision-maker needs. Selection is $\arg\max_a \hat\Delta_a(x)$.

\paragraph{Model.} Gradient-boosted decision
trees~\cite{friedman2001gbm,ke2017lightgbm}, one regressor per policy, with
hyperparameters fixed before evaluation and never tuned on a test fold. The
data are a few hundred context-level instances with full feedback from every
policy; a high-capacity model is neither needed nor defensible at this sample
size. We do not use graph neural networks, deep networks, reinforcement
learning or contextual bandits: the decision is a one-shot choice over four
alternatives with complete feedback, which none of those formalisms is required
for.

\paragraph{Baseline ladder.} B0 is the best fixed policy; B1 the mechanistic
rule; B2 a depth-3 decision tree; B3 multinomial logistic regression; B4 the
GBDT advantage selector; B5 the selective variant below. VBS is reported as a
reference. \textbf{The primary comparison is B4 against B1}, not B4 against
B0: B0 is a low bar, and beating it demonstrates only that the decision varies
with context at all. If the mechanistic rule wins, the finding is that the
decision structure is simple enough that a learner adds nothing --- and that is
reported as the result.

\paragraph{Leakage discipline.} Evaluation is leave-one-context-out. Inside
each fold, standardisation, the identity of the SBS, tree thresholds,
conformal quantiles and the support envelope are all recomputed from training
data alone. No seed of a held-out context appears in training. No realised
outcome of a test context is used as a feature; the eight descriptors are
functions of the context definition and cannot be computed from an outcome.

\paragraph{Abstention.} Two gates must both pass before the selector acts.
The \emph{support} gate measures the mean distance to the five nearest training
contexts in standardised descriptor space and abstains above the 95th
percentile of the training fold's own such distances. The \emph{confidence}
gate is a split-conformal interval on the predicted advantage
($\alpha=0.10$, calibrated on 30\% of the training fold); the selector acts only
when the lower bound on the advantage over the fixed policy is positive. If
either gate fails, the selector abstains and the fixed policy is used.

Split conformal prediction guarantees coverage under
\emph{exchangeability}~\cite{vovk2005alrw,lei2018conformal}. Covariate shift
breaks that assumption~\cite{shimodaira2000,tibshirani2019covshift}, so the
conformal gate is a calibrated instrument in-distribution and a heuristic
outside it. It is the support gate, not the conformal gate, that addresses
extrapolation, and neither detects concept shift. We state this as a limitation
of the design rather than presenting conformal prediction as a solution to
out-of-distribution behaviour.

\paragraph{Distribution shift.} Eight held-out splits are evaluated
separately, never pooled: demand above and below the training range, an
interior demand level as an interpolation control, an unseen information lag, an
unseen penetration, an unseen green ratio, training without disruption and
testing under it, and an incident located on a corridor that carries no
disruption in any training context. Pooling these would average an
interpolation control together with a domain shift and hide the only thing the
experiment can say about them, which is that they differ.
